\doublespacing % Do not change - required

\chapter{Literature review}
\label{ch2}

%%%%%%%%%%%%%%%%%%%%%%%%%%%%%%%%%%%%%%%
% IMPORTANT
\begin{spacing}{1} %THESE FOUR
\minitoc % LINES MUST APPEAR IN
\end{spacing} % EVERY
\thesisspacing % CHAPTER
% COPY THEM IN ANY NEW CHAPTER
%%%%%%%%%%%%%%%%%%%%%%%%%%%%%%%%%%%%%%%




\section{Fundamentals of LLM Quantization}

LLM quantization reduces the numerical precision used to represent model tensors in order to lower storage, memory traffic and potentially inference cost. The practical difficulty increases as precision is reduced, particularly for Transformer models whose weight and activation distributions may contain sensitive channels or large-magnitude outliers. Quantization can therefore differ not only in bit width but also in scope, ranging from weight-only schemes to methods that additionally quantize runtime activations and the \ac{KV} cache. These distinctions provide the basis for the specific methods reviewed in the following sections.


\subsection{Numerical precision and quantization}
Neural-network parameters and intermediate tensors are conventionally represented using floating-point formats such as FP32, FP16 or BF16. Lower-precision formats reduce the number of bits required for each value and can therefore decrease model storage and memory traffic. Quantization extends this principle by mapping a high-precision tensor onto a finite set of discrete reconstruction levels, commonly represented using integer formats such as INT8 or INT4. When appropriate low-precision kernels are available, this representation can additionally enable integer or mixed-precision computation rather than merely reducing checkpoint size \cite{jacob2017quantizationtrainingneuralnetworks}.

A commonly used uniform affine quantizer maps a real-valued element \(x\) to an integer code according to

\begin{equation}
q(x)
=
\operatorname{clip}
\left(
\operatorname{round}
\left(
\frac{x}{s}
\right)
+z,\,
q_{\min},q_{\max}
\right),
\label{eq:uniform-quantization}
\end{equation}

where \(s\) is the quantization scale, \(z\) is the zero-point, and \(q_{\min}\) and \(q_{\max}\) define the representable integer range. The approximate reconstructed value is then

\begin{equation}
\hat{x}
=
s\left(q(x)-z\right).
\label{eq:uniform-dequantization}
\end{equation}

For a \(b\)-bit representation, no more than \(2^b\) distinct integer codes are available. Reducing the bit width lowers storage requirements but also decreases the number of reconstruction levels available to approximate the original tensor. This introduces quantization error, and the error generally becomes more difficult to control as the precision is reduced. Quantization methods are consequently not only to approximate individual tensor values, but also to limit the effect of their approximation on the output and language-model behavior of the network.


\subsection{Post-training quantization}
\ac{PTQ} converts an already trained high-precision model into a
lower-precision representation without repeating the original
end-to-end training process. This distinguishes it from
quantization-aware training, in which simulated quantization effects
are introduced during optimisation so that the model parameters can
adapt to the resulting numerical error. Earlier integer-inference
frameworks demonstrated that incorporating quantization during training
can preserve model accuracy while enabling low-precision execution
\cite{jacob2017quantizationtrainingneuralnetworks}. However, repeating
such optimisation can be prohibitively expensive for large pretrained
language models, making \ac{PTQ} particularly attractive when the
original training data, training infrastructure or computational budget
are unavailable.

A typical \ac{PTQ} pipeline uses a comparatively small representative
calibration dataset to estimate the numerical ranges and statistical
properties of the tensors being quantized. These observations may be
used to determine scale factors, clipping thresholds, zero-points or
other quantizer parameters. Model weights are fixed after training and
can therefore be analysed and quantized offline. Activations, by
contrast, depend on the current input and may use either scales
estimated during calibration or dynamically calculated runtime scales.
The calibration data should be representative of the intended model
inputs, but must remain separate from the data used for final
evaluation.

\subsection{Scalar, group-wise and vector quantization}
Scalar and vector quantization differ primarily in the number of values
represented jointly, whereas group-wise quantization describes the
granularity over which quantization parameters are shared. In scalar
quantization, each tensor element is mapped independently to a discrete
reconstruction level. The affine quantizer introduced in
Equations~\eqref{eq:uniform-quantization}
and~\eqref{eq:uniform-dequantization} is an example of scalar
quantization: each real-valued element is represented by an integer code
together with the scale and zero-point required for reconstruction
\cite{jacob2017quantizationtrainingneuralnetworks}. This formulation is
computationally convenient and maps naturally onto integer and
mixed-precision hardware, but its accuracy depends strongly on how well
the selected quantization parameters represent the numerical range of
the tensor.

The quantization parameters may be defined at different granularities.
Per-tensor quantization uses a single scale and zero-point for an entire
tensor, resulting in low metadata overhead but limited adaptation to
local variations in magnitude. Per-channel quantization instead assigns
separate parameters to individual output or input channels. Group-wise
quantization provides an intermediate design by partitioning a tensor
into smaller groups, with each group using its own quantization
parameters. A smaller group size can represent local tensor ranges more
accurately, particularly when different regions of the tensor have
substantially different magnitudes. However, it also increases the
number of scales, zero-points or other auxiliary values that must be
stored and processed. Fine-grained group-wise quantization has therefore
been used in Transformer quantization systems such as ZeroQuant to
balance numerical accuracy, metadata overhead and hardware efficiency
\cite{yao2022zeroquantefficientaffordableposttraining}.

Vector quantization jointly represents a block of \(d\) values using a
single codeword from a multidimensional codebook. For an input vector
\(\mathbf{x}\in\mathbb{R}^{d}\) and a codebook
\(\mathcal{C}\subset\mathbb{R}^{d}\), nearest-neighbour vector
quantization may be expressed as

\begin{equation}
Q_{\mathrm{V}}(\mathbf{x})
=
\underset{\mathbf{c}\in\mathcal{C}}{\operatorname{arg\,min}}
\left\|
\mathbf{x}-\mathbf{c}
\right\|_{2}^{2},
\label{eq:vector-quantization}
\end{equation}

where the selected codeword \(Q_{\mathrm{V}}(\mathbf{x})\) provides the
reconstruction of the complete input vector\cite{1094577}. Unlike scalar
quantization, this representation can account for the joint geometry of
multiple tensor elements when allocating reconstruction points.
Consequently, vector quantization can offer a different
rate--distortion trade-off from independently quantizing each element.
Its practical use, however, requires storage or structured generation of
the codebook, encoding of codeword indices, and an execution path capable
of processing the resulting packed representation efficiently.

Lattice quantization is a structured form of vector quantization in
which the reconstruction points follow a regular geometric arrangement.
This structure can avoid storing an unrestricted learned codebook and
can provide efficient mathematical procedures for encoding and
reconstruction. NestQuant applies this principle using nested lattices
and low-dimensional vector blocks to quantize operands involved in
matrix products \cite{savkin2025nestquantnestedlatticequantization}. The distinction between scalar and vector quantization is therefore central to this dissertation:
GPTQ, AWQ, SmoothQuant and the project-specific W4A8 pipeline primarily
employ scalar or group-wise representations, whereas NestQuant
investigates a structured vector-quantization approach. The later
method-specific sections examine how these alternative representations
affect quantization error, effective storage and implementation
complexity.


\section{Weight-Only Post-Training Quantization}

Weight-only \ac{PTQ} reduces model storage and memory traffic by quantizing pretrained weights while retaining activations and the \ac{KV} cache at higher precision, with \acs{GPTQ} and \acs{AWQ} representing two influential approaches for limiting the resulting degradation in model quality.

\subsection{\ac{GPTQ}}
\ac{GPTQ} is a one-shot weight-only post-training quantization method
designed to preserve the behaviour of large Transformer models at low
weight precision. It extends the second-order quantization principles of
Optimal Brain Compression, which in turn adapts the Optimal Brain
Surgeon framework to post-training weight quantization
\cite{frantar2022obc,frantar2023gptq}. Rather than minimising the
element-wise difference between the original and quantized weights,
GPTQ quantizes each linear layer by attempting to preserve its output
over a small calibration dataset.

Let \(\mathbf{W}\) denote the full-precision weight matrix of a linear
layer and let \(\mathbf{X}\) contain representative input activations
collected during calibration. The layer-wise reconstruction problem can
be expressed as

\begin{equation}
\widehat{\mathbf{W}}
=
\underset{\widetilde{\mathbf{W}}\in\mathcal{Q}}
{\operatorname{arg\,min}}
\left\|
\mathbf{W}\mathbf{X}
-
\widetilde{\mathbf{W}}\mathbf{X}
\right\|_{F}^{2},
\label{eq:gptq-reconstruction}
\end{equation}

where \(\mathcal{Q}\) denotes the set of weight matrices permitted by
the selected quantization grid. This objective differs from directly
minimising
\(\|\mathbf{W}-\widetilde{\mathbf{W}}\|_{F}^{2}\), because it weights
the importance of a weight perturbation according to the calibration
activations passing through the layer. A numerically small weight error
can therefore still be important if it strongly affects the layer
output.

For each row of the weight matrix, the reconstruction objective is
quadratic and has an activation-dependent Hessian. GPTQ uses the
damped approximation

\begin{equation}
\mathbf{H}
=
2\mathbf{X}\mathbf{X}^{\top}
+
\lambda\mathbf{I},
\label{eq:gptq-hessian}
\end{equation}

where \(\lambda\mathbf{I}\) provides numerical damping. When a weight
or weight column is mapped to its quantized value, information from
\(\mathbf{H}^{-1}\) is used to update the remaining unquantized
weights. These updates compensate for the output error introduced by
the current rounding decision, allowing subsequent weights to absorb
part of the perturbation rather than quantizing all values
independently.

GPTQ introduces several modifications that make this second-order
procedure scalable to models containing billions of parameters. The
underlying Optimal Brain Quantization method greedily determines a
separate weight order for each row, which requires repeated and
expensive updates of the inverse Hessian. GPTQ instead quantizes all
rows according to the same fixed column order. Since the Hessian
depends on the calibration inputs rather than on the individual weight
rows, this permits the second-order information to be shared across
rows and substantially reduces computational complexity
\cite{frantar2023gptqaccurateposttrainingquantization}.

The method further processes consecutive columns in blocks and delays
updates to weights outside the current block until the block has been
quantized. These lazy batch updates convert many small memory-bound
operations into larger matrix operations that execute more efficiently
on a GPU. A Cholesky-based reformulation of the damped inverse Hessian
is additionally used to improve numerical stability when quantizing
large layers. Together, fixed-order processing, block-wise updates and
the Cholesky reformulation enable GPTQ to retain second-order error
compensation while scaling substantially beyond the model sizes
supported by earlier reconstruction-based quantizers.

GPTQ is compatible with different quantization grids and granularities.
For example, independent quantization parameters may be assigned to
groups of consecutive weights, allowing local weight ranges to be
represented more accurately than with a single grid for an entire row.
Smaller groups can improve low-bit reconstruction quality, particularly
at 3-bit and 4-bit precision, although they increase the number of
scales and other metadata stored alongside the packed weights.

The principal strength of GPTQ is therefore its ability to combine
layer-output-aware reconstruction with practical scalability. The
original study demonstrated accurate 3-bit and 4-bit weight
quantization for model families extending to hundreds of billions of
parameters \cite{frantar2023gptqaccurateposttrainingquantization}. However, GPTQ remains principally a weight-only method: runtime activations and the KV cache are not reduced by the core algorithm. Its output reconstruction also depends on
activation statistics collected from a calibration dataset, making the
representativeness of those samples relevant to the resulting
quantized model.

A further distinction must be made between compression and execution
speed. GPTQ can substantially reduce the amount of weight data stored
and transferred, but realised inference acceleration depends on the
availability of kernels that can load packed low-bit weights and
efficiently combine them with higher-precision activations. The
algorithm should therefore be evaluated both as a method for preserving
model quality after weight compression and as a representation whose
runtime benefit depends on the supporting execution backend.

\subsection{\ac{AWQ}}

\subsection{Critical comparison of weight-only approaches}


\section{Weight–Activation Quantization}

\subsection{Activation outliers}


\subsection{SmoothQuant}


\subsection{Low-bit \ac{W4A8}-related approaches}


\subsection{Limitations of joint W/A quantization}